\section{Exact observer-conditioned deal inference}
\label{sec:inference}

Every agent in this lineage is built on an inference object that this section defines formally,
because it is one of the paper's principal contributions and because its exactness is what makes the
policy results interpretable. The method is an exact enumeration carried out through a factorised
dynamic program. It is not an analytic closed-form expression, and the paper avoids that phrase.

\subsection{The engine, and the barrier between rules and policy}
\label{sec:inference-engine}

The engine is a single-header C++20 core. A hand is one 64-bit word over the 54 card indices
$c = 6h + i$ for half-suit $h$ and offset $i$, so hand arithmetic is bit arithmetic. Random streams
are counter-based and keyed by the triple (seed, seat, purpose), which gives the property everything
else depends on: a deal and its entire play reproduce from the seed alone. The duplicate-block design
of \S\ref{sec:evaluation-duplicate} rests on that property and on nothing else.

An acting policy never receives the game state. It receives its own knowledge object, its own hand
and the public event stream. There is no reference from a policy to the deal, to another seat's hand
or to the referee. This is an architectural barrier rather than a convention, and it is what allows
the side-channel audit of \S\ref{sec:evaluation-side} to be read as a test of the policy: a policy
that cannot reach the deal cannot read it, so any decision the audit fails to reproduce from public
information indicates state carried across deals rather than information taken from the referee.

\subsection{State, observations and constraints}
\label{sec:inference-constraints}

Let the hidden state be the initial deal $D$, an assignment of the 54 cards to the six seats with
nine cards each. Write $\mathcal{D}$ for the set of all such assignments; from one seat's viewpoint
at deal time $|\mathcal{D}| = 45!/(9!)^5 \approx 1.9 \times 10^{28}$.

An observer at seat $s$ holds an information set determined by five families of hard constraint,
each of which is a logical restriction on $D$ rather than a probabilistic one:

\begin{description}
\item[(C1) Own hand.] Seat $s$ knows its own nine cards exactly.
\item[(C2) Card counts.] The current number of cards at every seat is public, and given the transfer
  record this determines each seat's initial count.
\item[(C3) Transfers.] Every card that has changed hands, and between whom, is public.
\item[(C4) Failed asks.] For each failed ask, the target did not hold the named card at that moment,
  which given (C3) is a statement about $D$.
\item[(C5) Ask legality.] For each ask, the asker held at least one \emph{other} card of the named
  half-suit and did not hold the named card, again at that moment and hence about $D$.
\end{description}

The posterior of interest is uniform over the deals satisfying every constraint:
\[
  P(D \mid \mathcal{I}_s) \;=\; \frac{\mathbb{1}\!\left[D \models \mathcal{I}_s\right]}{Z},
  \qquad
  Z \;=\; \bigl|\{D \in \mathcal{D} : D \models \mathcal{I}_s\}\bigr| .
\]

\subsection{Factorisation by half-suit, joined by a count dynamic program}
\label{sec:inference-factor}

Computing $Z$ by direct enumeration is infeasible, and in general a constraint system of this shape
induces a permanent-like count, the object the literature approximates with inclusion-exclusion
\cite{ryser}, sign-vector formulae \cite{glynn}, Markov chains \cite{jsv}, matrix scaling \cite{lsw}
or sequential importance sampling \cite{chen-sis}.

Fish sidesteps all of it, because the constraints factorise by half-suit. Constraint (C5) is a
statement about a half-suit. Constraint (C4) is a statement about a card and therefore about a
half-suit. Constraint (C3) resolves within a half-suit. The \emph{only} coupling across half-suits is
(C2), the per-seat card counts, which is a linear constraint on how many cards each seat draws from
each block.

The computation therefore proceeds in two stages.

\begin{enumerate}
\item \textbf{Block enumeration.} For each half-suit $h$ independently, enumerate the assignments of
  its unresolved cards to seats that satisfy every (C1), (C3), (C4) and (C5) constraint touching
  $h$. Because a half-suit has six cards and resolved cards are removed, the surviving block is
  small. Accumulate, for each block, a table indexed by the vector of counts it contributes to each
  seat, recording how many assignments realise that count vector and, for each card and seat, how
  many of those place that card at that seat.
\item \textbf{Count join.} Combine the block tables by a dynamic program over the running count
  vector, which enforces (C2) exactly. The terminal cell of that program is $Z$; the same recursion
  carried with the per-card sub-counts yields the marginals.
\end{enumerate}

From this the engine reads off, exactly and without sampling: the normalisation constant $Z$; the
per-card marginal $P(\text{card } c \text{ at seat } j \mid \mathcal{I}_s)$; the probability that a
given team owns a whole half-suit; and the probability of a fully named allocation of a half-suit,
which is exactly the quantity a declaration decision requires.

An exact rebuild costs \vsixCountCost{}~ms, which is why the project has been able to treat inference
as solved and spend its cycles on policy.

\subsection{The deployed approximation and its policy prior}
\label{sec:inference-prior}

The exact posterior is exact \emph{under a uniform prior over constraint-consistent deals}. It
contains no model of behaviour, because the block construction consumes the deduction state alone
and has no parameter through which an opponent's tendencies could enter.

The path the agent actually deploys adds one. It applies a fitted soft reweighting of the rules
posterior with two parameters and no learned opponent model, in the spirit of the Skat inference
line \cite{solinas-supervised,rebstock-inference,buro-kermit} but far simpler; the deployed weight is $\theta = \vsixPriorTheta{}$. \S\ref{sec:background-vsix} reports the consequence, which is that the
approximation is the better predictor of card location, and \S\ref{sec:discussion-limits} discusses
what that implies.

\subsection{Verification}
\label{sec:inference-verify}

Three checks run over the engine, and all three ran as preconditions of the evaluation in
\S\ref{sec:results}.

\begin{description}
\item[Soundness audit.] Every seat's deduced knowledge is checked against the true deal at every
  decision, since no agent's deduction may ever exclude the truth. The verification run for this
  evaluation recorded \vsevenVerifyViolations{} violations in \vsevenVerifyChecks{} checks, and the
  preceding cycle \vsixAuditViolations{} in \vsixAuditChecks{}.
\item[Cross-engine agreement.] The block dynamic program, an independent card-counting path and an
  exact sampler must agree. The self-test reports \vsevenSelftestChecks{} checks passing, with
  block-versus-card agreement at the limit of double precision.
\item[Identity control.] The current policy, carrying its predecessor's parameter vector with every
  new mechanism disabled, must produce a transcript identical to its predecessor's under md5. It does
  (\vsixIdentity{}). Without this an ablation is not an ablation, because a measured difference could
  be an artefact of refactoring.
\end{description}

Two defects these checks found are recorded here because both changed previously published numbers,
and both survived for the same reason: they were invisible to the path the agent actually deploys.
The block dynamic program parked every instance's tables in a shared per-thread pool, so a second
construction on the same thread silently repointed the first instance's tables, giving
\numAliasMismatch{} mismatches in \numAliasChecks{} checks; any earlier number computed with the
exact belief in a configuration where more than one seat held such an object is affected, and the
deployed approximate-belief numbers are not. Separately, an audit that verifies the declaration
pre-gate never rejects a correct declaration had been parsed only in an older branch and so had never
run; once enabled it reports \vsixGateFalseNeg{} false negatives over \vsixGateOpportunities{}
declaration opportunities. Appendix~\ref{app:corrections} carries the full register.
